\begin{tikzpicture}[
    font=\small,
    >=stealth,
    block/.style={
        draw, 
        rounded corners=12pt, 
        minimum width=2.4cm, 
        minimum height=2.4cm, 
        align=center, 
        thick
    },
    maskblock/.style={
        draw, 
        rounded corners=12pt, 
        minimum width=2.4cm, 
        minimum height=2.4cm, 
        align=center, 
        thick
    }
]

    % ==========================================
    % 1. ENTRADA: IMAGEN RUIDOSA DE LA PASTILLA
    % ==========================================
    
    % Caja Imagen I (Ruidosa)
    \draw[thick, rounded corners=12pt] (-0.5, 0.25) rectangle (2.25, 2.93);
    \node[above, font=\bfseries] at (0.85, 3.6) {Imagen \(I\) (Ruidosa)};
    
    % Ilustración de la pastilla con ruido interno
    \draw[thick, fill=gray!20, dashed] (0.85, 1.6) circle (0.8cm);
    \node[font=\tiny, align=center] at (0.85, 1.6) {Pastilla\\+ Ruido};
    \node[font=\tiny, gray] at (0, 2.73) {\(\cdot \cdot \times\)};
    \node[font=\tiny, gray] at (1.73, 0.72) {\(\times \cdot \cdot\)};

    % Flecha hacia Banco de Gabor / Matriz X
    \draw[thick, ->] (2.3, 1.6) -- node[above, font=\tiny] {16 Filtros} node[below, font=\tiny] {Gabor} (3.2, 1.6);


    % ==========================================
    % 2. MATRIZ X Y K-MEANS (TEXTONES)
    % ==========================================

    % Matriz X de Respuestas
    \draw[thick, rounded corners=8pt] (3.3, -0.4) rectangle (4.9, 3.6);
    \node[above, font=\bfseries] at (4.1, 3.6) {Matriz \(X\)};
    
    \node at (4.1, 3.0) {\tiny \([r_1 \dots r_{16}]_0\)};
    \node at (4.1, 2.3) {\tiny \([r_1 \dots r_{16}]_1\)};
    \node at (4.1, 1.6) {\(\vdots\)};
    \node at (4.1, 0.9) {\tiny \([r_1 \dots r_{16}]_j\)};
    \node at (4.1, 0.2) {\(\vdots\)};
    \node at (4.1, -0.2) {\tiny \([r_1 \dots r_{16}]_N\)};

    % Flecha hacia KMeans
    \draw[thick, ->] (4.9, 1.6) -- (5.7, 1.6);

    % Bloque KMeans Textons
    \node[block, minimum width=2.4cm] (km1) at (6.9, 1.6) {
        \textbf{K-Means}\\[0.1cm]
        \tiny Textones\\
        \tiny (\(K_1 = 10\))
    };

    % Flecha a Mapa de Textones
    \draw[thick, ->] (km1) -- (8.46, 1.6);


    % ==========================================
    % 3. MAPA DE TEXTONES Y HISTOGRAMAS
    % ==========================================

    % Mapa de Textones (H x W)
    \draw[thick, rounded corners=8pt] (8.5, -0.4) rectangle (10.1, 3.6);
    \node[anchor=north, font=\bfseries, align=center] at (9.3, 3.5) {\footnotesize Texton\\[-0.05cm]\footnotesize Map};
    
    % Ventana 21x21 sobre el Texton Map
    \draw[dashed, red!80!black, thick] (8.8, 1.1) rectangle (9.8, 2.1);
    \node[font=\tiny, text=red!80!black] at (9.3, 0.8) {Ventana \(21\times21\)};

    % Flecha hacia Histograma Local
    \draw[thick, ->] (10.2, 1.6) -- (11.0, 1.6);

    % Matriz de Histogramas Locales
    \draw[thick, rounded corners=8pt] (11.1, -0.4) rectangle (12.8, 3.6);
    \node[anchor=north, font=\bfseries, align=center] at (11.95, 3.5) {\footnotesize Matriz\\[-0.05cm]\footnotesize Hist.};
    
    \node at (11.95, 2.4) {\tiny \(\mathbf{h}_0 \in \mathbb{R}^{10}\)};
    \node at (11.95, 1.8) {\tiny \(\mathbf{h}_1 \in \mathbb{R}^{10}\)};
    \node at (11.95, 1.3) {\(\vdots\)};
    \node at (11.95, 0.75) {\tiny \(\mathbf{h}_j \in \mathbb{R}^{10}\)};
    \node at (11.95, 0.15) {\(\vdots\)};
    \node at (11.95, -0.15) {\tiny \(\mathbf{h}_N \in \mathbb{R}^{10}\)};

    % Flecha hacia GMM
    \draw[thick, ->] (12.9, 1.6) -- (13.7, 1.6);


    % ==========================================
    % 4. GMM Y SALIDA MÁSCARA (TU BLOQUE ORIGINAL)
    % ==========================================

    % Bloque GMM
    \node[block, minimum width=2.4cm] (gmm) at (14.9, 1.6) {
        \textbf{GMM}\\[0.1cm]
        \tiny Regiones\\
        \tiny (\(K_2 = 10\))
    };

    % Bloque Mask (Misma construcción e identidad visual exacta que el anterior)
    \node[maskblock, right=1.02cm of gmm, minimum width=78.4613287pt, minimum height=79.6677171pt] (mask) {};
    \node[anchor=north, font=\bfseries] at (mask.north) {Mask};

    % Coordenada auxiliar para el centro del dibujo de la máscara
    \coordinate (mc) at ([yshift=-0.25cm]mask.center);

    % Dibujo interno en Mask
    \draw[dashed, thick] ([xshift=-0.45cm, yshift=-0.4cm]mc) rectangle ([xshift=0.45cm, yshift=0.4cm]mc);
    \filldraw[fill=green!50!white, draw=black, thick] (mc) circle (0.9cm);
    
    % Onda blanca interna
    \draw[white, line width=1.2pt] ([xshift=-0.2cm]mc) sin ([xshift=-0.1cm, yshift=0.08cm]mc) cos (mc);
    \draw[white, line width=1.2pt] (mc) sin ([xshift=0.1cm, yshift=-0.08cm]mc) cos ([xshift=0.2cm]mc);

    % Conexión final desde GMM hacia Mask
    \draw[thick, ->] (gmm) -- (mask.west);

\end{tikzpicture}